\section{Integração}

A integração dos componentes, conforme mencionado anteriormente, ocorre por meio da agregação de todos os recursos de que uma \gls{vm} necessita em sua própria classe.
Além disso, cada mediador virtual possui um mapeamento específico para um mediador real, como ilustrado na Figura~\ref{fig:virtual-to-real}.

\begin{figure}[H]
\centering
\includegraphics[width=0.6\textwidth]{images/VirtualToReal.pdf}
\caption{Diagrama de Mapeamento de Componentes Virtuais para Reais}\label{fig:virtual-to-real}
\end{figure}

Formalizando essa estrutura, a interface mínima que a camada de virtualização fornece é apresentada na Figura~\ref{fig:virtualization-framework-uml}.

\begin{figure}[H]
\centering
\includegraphics[width=0.8\textwidth]{images/VirtualizationFrameworkUML.pdf}
\caption{Diagrama de Classes Relacionadas ao Processo de Virtualização}\label{fig:virtualization-framework-uml}
\end{figure}
